% Image credits for the photographs and AI illustrations of Book 5
% (University Biology, Year 3). Machine-readable license records live in
% images/book5/CREDITS.md; keep the two files in sync.
\chapter*{Beeldverantwoording}

% Under bidi=basic a Latin run is ONE directional box, so a long author or
% institution name cannot be broken at all; \allowbreak inside it does nothing.
% The Arabic edition reported this page as the book's only overfull boxes.
% \sloppy must be in force when the paragraph is BROKEN, hence the \par before
% \endgroup -- without it the group closes first and the tolerance is restored.
\begingroup\sloppy

De tekeningen in dit boek zijn origineel. De foto's worden gebruikt
onder hun respectieve vrije licenties, met dank aan hun makers:

\begin{itemize}
  \item Hoofdstuk~\ref{ch:b3:genomics}: foto van Frederick Sanger ---
        publiek domein (Wikimedia Commons).
  \item Hoofdstuk~\ref{ch:b3:genetic-engineering}: gouden rijst naast
        gewone rijst, International Rice Research Institute, 2011 ---
        CC~BY~2.0 (Wikimedia Commons).
  \item Hoofdstuk~\ref{ch:b3:structural-biology}: foto van Max Perutz,
        Associated Press, 1962 --- publiek domein (Wikimedia Commons).
  \item Hoofdstuk~\ref{ch:b3:bacteriology}: portret van Robert Koch,
        Wellcome Collection --- CC~BY~4.0 (Wikimedia Commons); foto van
        Alexander Fleming in zijn laboratorium, archief van het
        U.S. Navy Bureau of Medicine and Surgery --- geen bekende
        auteursrechtelijke beperkingen (Wikimedia Commons).
  \item Hoofdstuk~\ref{ch:b3:innate-immunity}: foto van Élie
        Metchnikoff, Bain News Service, Library of Congress --- publiek
        domein (Wikimedia Commons).
  \item Hoofdstuk~\ref{ch:b3:adaptive-immunity}: gegraveerd portret van
        Edward Jenner, 1807 --- publiek domein (Wikimedia Commons).
  \item Hoofdstuk~\ref{ch:b3:nervous-systems}: tekening van een
        purkinjecel door Santiago Ramón y Cajal, 1899 --- publiek
        domein (Wikimedia Commons).
  \item Hoofdstuk~\ref{ch:b3:endocrinology}: foto van Frederick Banting
        en Charles Best, ca.~1924 --- publiek domein (Wikimedia
        Commons).
  \item Hoofdstuk~\ref{ch:b3:molecular-evolution}: foto van Motoo
        Kimura, 1986, fotograaf onbekend, via S.~Kumar en T.~Gojobori,
        \emph{Molecular Biology and Evolution} (2023) --- CC~BY~4.0
        (Wikimedia Commons).
  \item Hoofdstuk~\ref{ch:b3:behavioural-ecology}: foto van Niko
        Tinbergen door Rob Mieremet, Anefo, 1973 --- CC0, Nationaal
        Archief (Wikimedia Commons).
  \item Hoofdstuk~\ref{ch:b3:conservation-biology}: foto van Martha, de
        laatste trekduif, Enno Meyer, 1912 --- publiek domein (Wikimedia
        Commons); wolven in Yellowstone National Park, National Park
        Service, 1996 --- publiek domein (Wikimedia Commons); Sirocco de
        kakapo, Department of Conservation, Nieuw-Zeeland, 2012 ---
        CC~BY~2.0 (Wikimedia Commons).
\end{itemize}

De volledige bronvermeldingen en licentieadressen van elke foto staan
in \texttt{images/\allowbreak book5/\allowbreak CREDITS.md} in de repository van het boek.

\bigskip
Naast de foto's en de originele tekeningen zijn ongeveer negentig
figuren verspreid over de hoofdstukken van dit deel illustraties die
met kunstmatige intelligentie zijn gemaakt, met de beeldgeneratie van
OpenAI en op basis van prompts die de redacteuren van het boek hebben
geschreven, en die vóór opname op biologische juistheid zijn
nagekeken. De volledige prompt van elk gegenereerd beeld staat in
\texttt{images/\allowbreak book5/\allowbreak ai/\allowbreak PROMPTS.md} in de repository.
\par\endgroup
\clearpage
